\documentclass[11pt,a4paper]{article}

\usepackage[utf8]{inputenc}
\usepackage[margin=0.75in]{geometry}
\usepackage{xcolor}
\usepackage{listings}
\usepackage{fancyhdr}
\usepackage{hyperref}
\usepackage{titlesec}

% Professional Color Palette
\definecolor{codegreen}{rgb}{0.12, 0.55, 0.25}
\definecolor{codegray}{rgb}{0.45, 0.45, 0.45}
\definecolor{codepurple}{rgb}{0.55, 0.15, 0.70}
\definecolor{codeblue}{rgb}{0.10, 0.35, 0.75}
\definecolor{backcolour}{rgb}{0.97, 0.98, 0.99}
\definecolor{rulecolor}{rgb}{0.80, 0.85, 0.90}

% Global Listings Settings
\lstdefinestyle{customcode}{
    backgroundcolor=\color{backcolour},
    commentstyle=\color{codegreen}\itshape,
    keywordstyle=\color{codeblue}\bfseries,
    numberstyle=\tiny\color{codegray},
    stringstyle=\color{codepurple},
    basicstyle=\ttfamily\footnotesize,
    breakatwhitespace=false,
    breaklines=true,
    captionpos=t,
    keepspaces=true,
    numbers=left,
    numbersep=8pt,
    showspaces=false,
    showstringspaces=false,
    showtabs=false,
    tabsize=4,
    frame=single,
    rulecolor=\color{rulecolor},
    framerule=0.8pt,
    xleftmargin=15pt,
    framexleftmargin=15pt
}

\lstset{style=customcode}

% Page Header/Footer
\pagestyle{fancy}
\fancyhf{}
\lhead{\textbf{ESP32 IoT Weather Monitoring \& ML Forecasting System}}
\rhead{Source Code Appendix}
\cfoot{\thepage}

\hypersetup{
    colorlinks=true,
    linkcolor=codeblue,
    urlcolor=codepurple,
    pdftitle={ESP32 Weather Prediction - Complete Source Code Appendix},
    pdfauthor={Group 11 - NSS College of Engineering, Palakkad}
}

\begin{document}

\begin{titlepage}
    \centering
    \vspace*{2cm}
    
    {\Huge \textbf{ESP32 IoT Weather Monitoring \&\\ AI Prediction System}}\\[0.8cm]
    {\LARGE \textbf{Complete Source Code Appendix}}\\[1.5cm]
    
    \textbf{Department of Electrical and Electronics Engineering}\\[0.3cm]
    \textbf{NSS College of Engineering, Palakkad}\\[1.5cm]
    
    \textbf{Group 11}\\[0.4cm]
    \textbf{Rahul Das V R} (Reg No: NSS23EE107)\\[2.5cm]
    
    \vfill
    {\large Multi-Horizon Environmental Forecasting: 10, 20, and 30-Minute Horizons}\\[0.2cm]
    {\large Powered by Multi-Output XGBoost and Deep Sequence LSTM Networks}\\[0.5cm]
    \today
\end{titlepage}

\tableofcontents
\newpage

\section{System Architecture \& Repository Layout}
This document serves as the comprehensive source code appendix for the \textbf{ESP32 IoT Weather Monitoring and Predictive Intelligence System}. The codebase is structured into five core functional modules:

\begin{enumerate}
    \item \textbf{Data Acquisition \& Telemetry:} Continuous local logging from ESP32 via HTTP REST API.
    \item \textbf{Data Sanitization \& Preprocessing:} Glitch rejection, multi-horizon target formation.
    \item \textbf{Model Training Pipelines:} Boosted tree regression (XGBoost) and sequential neural networks (LSTM).
    \item \textbf{Production Server \& REST API:} Flask application hosting live telemetry, background logging, and real-time model inference.
    \item \textbf{Front-End Speedometer Dashboard:} Cockpit-style SVG speedometer dials, animated needles, and real-time Chart.js tracking.
\end{enumerate}

\newpage
\section{Data Pipeline Modules}

\subsection{1. \texttt{clean\_data.py} -- Data Sanitization and Glitch Filtering}
Validates sensor bounds (temperature between $-40^\circ\text{C}$ and $+85^\circ\text{C}$, humidity within $0-100\%$, plausible atmospheric pressures), resolves missing values, and removes duplicate records.

\lstinputlisting[language=Python, caption={clean\_data.py}]{clean_data.py}

\newpage
\subsection{2. \texttt{prepare\_data.py} -- Multi-Horizon Training Dataset Generator}
Pairs each historical reading with sensor observations 10, 20, and 30 minutes in the future using asynchronous nearest-neighbor temporal joins (\texttt{pd.merge\_asof}).

\lstinputlisting[language=Python, caption={prepare\_data.py}]{prepare_data.py}

\newpage
\subsection{3. \texttt{data\_logger.py} -- Standalone ESP32 Telemetry Logger}
Polls the ESP32 endpoint at \texttt{http://10.163.181.58/api/data} every 30 seconds and writes verified records to \texttt{weather\_data.csv}.

\lstinputlisting[language=Python, caption={data\_logger.py}]{data_logger.py}

\newpage
\section{Machine Learning Training Pipelines}

\subsection{1. \texttt{train\_xgboost.py} -- Multi-Horizon XGBoost Regressors}
Trains multi-output tree ensembles for 10-minute, 20-minute, and 30-minute horizons using temperature, humidity, pressure, and cyclical diurnal time features ($\sin/\cos$ encoding).

\lstinputlisting[language=Python, caption={train\_xgboost.py}]{train_xgboost.py}

\newpage
\subsection{2. \texttt{train\_lstm.py} -- Multi-Horizon Sequential LSTM Networks}
Normalizes multi-variate time series windows ($W = 10$ steps) through \texttt{MinMaxScaler} and trains recurrent memory cells to capture nonlinear physical weather trajectories.

\lstinputlisting[language=Python, caption={train\_lstm.py}]{train_lstm.py}

\newpage
\section{Production Backend \& Inference Server}

\subsection{1. \texttt{app.py} -- Flask API, Background Logger, \& Multi-Horizon Engine}
Runs the production application server on port 5001. Incorporates thread-safe deduplication logging, sensor telemetry caching, Open-Meteo external benchmarks, and real-time multi-model evaluation.

\lstinputlisting[language=Python, caption={app.py}]{app.py}

\newpage
\section{Front-End Speedometer Dashboard}

\subsection{1. \texttt{dashboard.html} -- Interactive Cockpit Speedometer Gauges}
The single-page user interface featuring automotive-grade SVG speedometer dials with rotating needle mechanics, live ESP32 status updates, and 6 separate multi-horizon forecast cards.

\lstinputlisting[language=html, caption={dashboard.html}]{dashboard.html}

\end{document}
